% !TeX root = ../xdupgthesis_template_lc.tex

\section*{1.\quad 基本情况}

\section*{2.\quad 教育背景}

\section*{3.\quad 攻读硕士学位期间的研究成果}

\subsection*{3.1\quad 申请（受理）专利及软件著作权}

[1] XXX 等. 一种低功耗地磁车辆检测方法及系统：中国，申请号：2026102390620，申请日期：2026 年 02 月 28 日（共 6 位发明人，本人为第二发明人，学生排名第一）。

[2] XXX 等. 面向多次漏检的双车道外侧地磁车辆轨迹连续重构方法及系统：中国，申请号：2026102390508，申请日期：2026 年 02 月 28 日（共 6 位发明人，本人为第三发明人，学生排名第二）。

[3] XXX 等. 单地磁异常停车检测与事件上报软件 V1.0：中国，受理号：2026R11S0463698，申请日期：2026 年 02 月 03 日（共 5 位作者，本人排名第二，学生排名第一）。

[4] XXX 等. 基于地磁传感器的拥堵车辆检测仿真系统 V1.0：中国，受理号：2026R11S0463768，申请日期：2026 年 02 月 03 日（共 6 位作者，本人排名第三，学生排名第二）。

[5] XXX 等. 基于深度学习的小车换道意图识别和轨迹预测系统 V1.0：中国，受理号：2026R11S0463705，申请日期：2026 年 02 月 04 日（共 7 位作者，本人排名第四）。

\subsection*{3.2\quad 参与科研项目及获奖}

[1] 国家重点研发计划项目，高速公路智能车路协同系统集成应用（SQ2019YFB160057），2020.09--2022.12，已结题，负责内容待补充。

[2] 国家自然科学基金重点支持项目，车联网环境的智能汽车多模态高精度融合定位（U21A20446），2022.01--2025.12，项目参与内容待补充。

[3] 科研获奖与奖学金信息待补充。


